\documentclass[11pt]{article}
\usepackage{amsmath,amssymb,amsthm,mathtools}
\usepackage[margin=1in]{geometry}
\usepackage[colorlinks=true,linkcolor=blue,citecolor=blue,urlcolor=blue]{hyperref}

\newtheorem{theorem}{Theorem}[section]
\newtheorem{lemma}[theorem]{Lemma}
\newtheorem{corollary}[theorem]{Corollary}
\newtheorem{proposition}[theorem]{Proposition}
\theoremstyle{definition}
\newtheorem{definition}[theorem]{Definition}
\theoremstyle{remark}
\newtheorem{remark}[theorem]{Remark}

\newcommand{\CC}{\mathbb{C}}
\newcommand{\RR}{\mathbb{R}}
\newcommand{\CP}{\mathbb{CP}}
\newcommand{\Hom}{\operatorname{Hom}}
\newcommand{\dd}{\mathrm{d}}

\begin{document}

\title{Problem 4 -- Humanize}
\author{}
\date{}
\maketitle

\begin{abstract}
\noindent We prove that the answer to the question is \emph{no} in
general. If $(M,\omega)$ is a symplectic manifold, $D\subset M$ a
symplectic simple normal crossings divisor, and $J$ an $\omega$-tame
almost complex structure admitting, locally on $M$, smooth complex
coordinates $(z_1,\ldots,z_n)$ with $z_1,\ldots,z_r$
$J$-holomorphic and $D=\{z_1\cdots z_r=0\}$, then the normal bundle
of every compact connected component $D_\alpha$ of $D$ admits a flat
structure with locally constant (locally) $J$-holomorphic trivialising
frames; in particular
\[
\langle e(\nu_{D_\alpha}),[D_\alpha]\rangle\ =\
D_\alpha\cdot D_\alpha\ =\ 0 .
\]
For a projective line $L\cong\CP^1$ in $(\CP^2,\omega_{\mathrm{FS}})$
(a symplectic SNC divisor with one component) one has $L\cdot L=1$,
so no such $J$ exists, and the cohomologous deformation of $\omega$
allowed in the question cannot remove the obstruction, which is
topological in the pair $(M,D)$.
\end{abstract}

\section{Statement and answer}\label{sec:statement}

\begin{quote}
\textbf{Question.} Let $(M,\omega)$ be a closed symplectic manifold
and $D\subset M$ a symplectic simple normal crossings divisor. Does
there exist an $\omega$-tame almost complex structure $J$ such that
locally on $M$ there are smooth complex coordinates $(z_1,\ldots,z_n)$
in which, for some $0\le r\le n$, the functions $z_1,\ldots,z_r$ are
$J$-holomorphic and $D$ is given by $z_1\cdots z_r=0$? (One may first
replace $\omega$ by a symplectic form connected to it through a path
of symplectic forms, all cohomologous to $\omega$, along which $D$
remains a symplectic simple normal crossings divisor.)
\end{quote}

\medskip
\noindent\textbf{Answer: No in general.} The precise form is:

\begin{theorem}[Obstruction]\label{thm:obstruction}
Let $(M,\omega)$ be a symplectic manifold, $D\subset M$ a symplectic
simple normal crossings divisor, and let $J$ be an almost complex
structure on $M$ which is \emph{SNC-adapted to} $D$, meaning: every
point $p\in M$ has a neighbourhood $U$ with smooth complex
coordinates $(z_1,\ldots,z_n)\colon U\to U'\subset\CC^n$ such that,
for some $0\le r\le n$,
\begin{enumerate}
\item[\normalfont(i)] $D\cap U=\{q\in U:z_1(q)\cdots z_r(q)=0\}$;
\item[\normalfont(ii)] $z_1,\dots,z_r$ are $J$-holomorphic on $U$
(the condition $\dd z_k\circ J=i\,\dd z_k$ holds at every point of
$U$).
\end{enumerate}
Let $D_\alpha$ be a compact connected component of $D$. Then the
normal bundle $\nu_\alpha$ of $D_\alpha$ in $M$ --- a complex line
bundle via the quotient of $J$, see Lemma~\ref{lem:strata} ---
possesses a flat connection; in particular its Euler class satisfies
\[
e(\nu_\alpha)\ =\ 0\ \in\ H^2(D_\alpha;\RR),
\qquad\text{hence}\qquad
\langle e(\nu_\alpha),[D_\alpha]\rangle
\ =\ D_\alpha\cdot D_\alpha\ =\ 0 .
\]
(The conclusion of the question would therefore force
$D_\alpha\cdot D_\alpha=0$ for every compact component, regardless
of the tameness of $J$ --- the latter enters the question but the
obstruction needs no taming at all.)
\end{theorem}

\begin{corollary}[The counterexample]\label{cor:cp2}
Let $M=\CP^2$, $\omega=\omega_{\mathrm{FS}}$ the Fubini--Study
symplectic form, and $D=L\cong\CP^1$ a projective line. Then $D$ is
a symplectic simple normal crossings divisor in the closed symplectic
manifold $(M,\omega)$: it is a smooth symplectic submanifold, and a
single component trivially constitutes simple normal crossings. But
\[
L\cdot L\ =\ \int_L c_1\bigl(\mathcal O_{\CP^1}(1)\bigr)\ =\ 1
\ \neq\ 0 ,
\]
since $\nu_L\cong\mathcal O_{\CP^1}(1)$. By
Theorem~\ref{thm:obstruction}, no $\omega'$-tame SNC-adapted $J$
exists for \emph{any} symplectic form $\omega'$ in a cohomologous
deformation family: the quantity $\langle e(\nu_L),[L]\rangle=L\cdot
L$ is a topological intersection number of the pair $(M,L)$, hence
unchanged by any deformation of $\omega$ that keeps $L$ smoothly
embedded (and a fortiori symplectic SNC). Moreover, the same
conclusion holds for the nodal SNC divisor $D=L_1\cup L_2$ formed by
two distinct lines: each component satisfies
$L_i\cdot L_i=1$.
\end{corollary}

\begin{remark}[What happens for components with trivial normal Euler
class]
For components with $e(\nu_\alpha)=0$ in $H^2(D_\alpha;\RR)$ (e.g.\
trivial normal bundle), Theorem~\ref{thm:obstruction} says nothing;
and indeed in the product situation $M=X\times\CC$ with
$D=X\times\{0\}$ an $\omega$-tame adapted $J$ does exist whenever
$(X,\omega|_X)$ carries a tame almost complex structure $J_X$: take
$J=J_X\oplus i$. This shows the obstruction found here is specific
to components with nontrivial normal Euler class, and the
counterexample relies on $L\cdot L\neq0$.
\end{remark}

\section{Strata are $J$-invariant}\label{sec:strata}

\begin{lemma}\label{lem:strata}
Let $J$ be SNC-adapted to $D$ (no taming assumption needed). Then
for every $p\in D$, and $A:=\{\alpha:p\in D_\alpha\}$, the tangent
space $T_pD_A$ of the stratum $D_A=\bigcap_{\alpha\in A}D_\alpha$
is a $J_p$-complex subspace of $T_pM$. Consequently
$J_A:=J|_{TD_A}$ is an almost complex structure on $D_A$, and, if
$J$ is $\omega$-tame, so is $J_A$ (with respect to
$\omega|_{D_A}$).
\end{lemma}

\begin{proof}
Choose adapted coordinates $(z_1,\dots,z_n)$ at $p$ with
$D\cap U=\{z_1\cdots z_r=0\}$, so that $D_A\cap U=\{z_{k}=
0:\alpha_k\in A\}$ after numbering, and
$T_pD_A=\bigcap_{k:\,\alpha_k\in A}\ker\dd z_k(p)$. For
$v\in\ker\dd z_k(p)$, the $J$-holomorphicity of $z_k$ gives
$\dd z_k(Jv)=i\,\dd z_k(v)=0$, so $Jv\in\ker\dd z_k(p)$. Thus each
$\ker\dd z_k(p)$ is $J_p$-invariant, and so is their intersection.
The statement about taming restricts from $M$ to the invariant
subspace: for $v\in T_pD_A\setminus0$,
$\omega(v,J_Av)=\omega(v,Jv)>0$.
\end{proof}

\begin{corollary}[Normal bundle is a complex line bundle]
\label{cor:normal}
For any component $D_\alpha$, the normal bundle
$\nu_\alpha:=TM|_{D_\alpha}/TD_\alpha$ carries a fibrewise complex
structure induced by $J$, making it a complex line bundle over
$D_\alpha$; the conormal bundle $\nu_\alpha^*=\Hom_\CC(\nu_\alpha,
\CC)$ is the dual complex line bundle.
\end{corollary}

\section{The ratio lemma}\label{sec:ratio}

\begin{lemma}[Ratios are $J$-holomorphic across the divisor]
\label{lem:ratio}
Let $J$ be SNC-adapted, and let $\zeta,\zeta'$ be $J$-holomorphic
functions on a connected open set $W\subset M$ with
\[
\{\zeta=0\}\ =\ \{\zeta'=0\}\ =\ W\cap D_\alpha,
\qquad \dd\zeta\neq0,\ \dd\zeta'\neq0\ \text{on }W\cap D_\alpha .
\]
Then the ratio $h:=\zeta'/\zeta$, smooth and nonvanishing on
$W\setminus D_\alpha$, extends to a smooth nonvanishing function on
$W$, the extension satisfies
\begin{equation}\label{eq:cross}
\dd\zeta'\ =\ h\,\dd\zeta\quad\text{on }W\cap D_\alpha,
\end{equation}
and $h$ is $J$-holomorphic on $W$: $\dd h\circ J=i\,\dd h$.
\end{lemma}

\begin{proof}
Off $D_\alpha$, both functions are nowhere zero, so $h$ is smooth
nonvanishing there, and the computation
$\dd\zeta'=\dd h\,\zeta+h\,\dd\zeta$ composed with $J$ gives, using
$\dd\zeta\circ J=i\dd\zeta$ and $\dd\zeta'\circ J=i\dd\zeta'$,
\[
0\ =\ \zeta\cdot\bigl(\dd h\circ J-i\,\dd h\bigr),
\]
hence $\dd h\circ J=i\dd h$ on $W\setminus D_\alpha$: $h$ is
$J$-holomorphic there.

For the extension: $D_\alpha\cap W$ is a smooth codimension-two
submanifold along which both $\zeta,\zeta'$ vanish with nonvanishing
differential transversely. In any tubular neighbourhood with normal
coordinate, both ${\zeta}$ and ${\zeta}'$ have a simple zero along
the zero section; their quotient extends smoothly and nonvanishingly
across the zero section (one-variable Taylor argument in the normal
direction: if $f(t)=t\,F(t)$, $g(t)=t\,G(t)$ with $F(0),G(0)\neq0$,
then $f/g=F/G$ extends). Differentiating
$\zeta'=h\zeta$ at $q\in D_\alpha$ (where $\zeta=0$) yields
\eqref{eq:cross}, and since $\dd\zeta'(q)\neq0$ and
$\dd\zeta(q)\neq0$, the extended $h$ is nonvanishing also on
$D_\alpha$.

Finally, the identity $\dd h\circ J=i\dd h$ holds on
$W\setminus D_\alpha$, both sides are continuous on $W$ (since $h$
and $J$ are smooth on $W$), and $W\setminus D_\alpha$ is dense in
$W$; hence the identity holds on all of $W$.
\end{proof}

\section{The maximum principle}\label{sec:max}

\begin{lemma}\label{lem:max}
Let $(X,J)$ be a connected almost complex manifold of dimension
$2m\ge2$, and $h\colon X\to\CC$ $J$-holomorphic. If $|h|$ has a
local maximum at some point of $X$, then $h$ is constant. In
particular, on a compact connected almost complex manifold every
$J$-holomorphic function is constant.
\end{lemma}

\begin{proof}
We use the local existence of $J$-holomorphic discs in almost
complex manifolds (Nijenhuis--Woolf): for every $p\in X$ and every
complex line $\ell\subset(T_pX,J_p)$ there is a $J$-holomorphic disc
$\varphi\colon\Delta\to X$ with $\varphi(0)=p$ and
$\varphi_{*}T_0\Delta=\ell$; moreover one can choose such discs
smoothly depending on parameters (e.g.\ on the line and on nearby
points), so that their images through $\{p'\}\times\{$lines in
$T_{p'}X\}$ cover a neighbourhood of $p$.

Suppose $|h|$ attains a local maximum at $p$. For any small
$J$-holomorphic disc $\varphi\colon\Delta\to X$ with
$\varphi(0)=p$, the composition $h\circ\varphi$ is holomorphic
$\Delta\to\CC$ (chain rule from $\dd h\circ J=i\dd h$), and
$|h\circ\varphi|$ attains a maximum at $0$; by the classical
maximum modulus principle $h\circ\varphi$ is constant, in
particular $h|_{\varphi(\Delta)}\equiv h(p)$. Since such discs
through $p$ sweep a neighbourhood of $p$ (choose the disc family as
above), $h\equiv h(p)$ on an open set containing $p$. Thus
$\{q:h(q)=h(p)\}$ is open; it is also closed, and $X$ is
connected, so $h\equiv h(p)$. The compact case follows because
$|h|$ attains its maximum.
\end{proof}

\section{Proof of the obstruction}\label{sec:proof}

\begin{proof}[Proof of Theorem~\ref{thm:obstruction}]
Let $D_\alpha$ be a compact connected component; all components are
smooth codimension-$2$ symplectic submanifolds, and we use the
conormal complex line bundle $\nu_\alpha^*$ of
Corollary~\ref{cor:normal}.

\emph{Step 1: $J$-holomorphic conormal frames.}
By SNC-adaptedness, cover $D_\alpha$ by finitely many adapted
coordinate neighbourhoods $W_i$ with the property that
$z_1^{(i)}$ is $J$-holomorphic and
$\{z_1^{(i)}=0\}=D_\alpha\cap W_i$: at any $p\in D_\alpha$ the
component $D_\alpha$ is among the coordinate hyperplanes of every
adapted coordinate system at $p$ (it appears among the
$D=\{z_1\cdots z_r=0\}$ factors), so renumbering puts it first;
this works also at points of $D_\alpha$ lying on other components.
On $W_i\cap D_\alpha$, set
\[
e_i\ :=\ \dd z_1^{(i)}\big|_{D_\alpha\cap W_i}
\ \in\ \Gamma\bigl(D_\alpha\cap W_i,\ T^*M\otimes\CC\bigr).
\]
Since $z_1^{(i)}|_{D_\alpha\cap W_i}=0$, $e_i$ annihilates
$T(D_\alpha)$, so it descends to a homomorphism
$\nu_\alpha\to\CC$; since $\dd z_1^{(i)}\neq0$ along $D_\alpha$
(it is a coordinate differential), $e_i\neq0$ everywhere on
$D_\alpha\cap W_i$; and by $J$-holomorphicity,
$e_i(Jv)=\dd z_1^{(i)}(Jv)=i\,\dd z_1^{(i)}(v)=i\,e_i(v)$, so $e_i$
is $\CC$-linear for the complex structure of
Corollary~\ref{cor:normal}: $e_i$ is a nowhere vanishing local frame
of the complex line bundle $\nu_\alpha^*$.

\emph{Step 2: constant transition functions.}
On $W_i\cap W_j$, both $z_1^{(i)}$ and $z_1^{(j)}$ are
$J$-holomorphic defining functions of
$D_\alpha\cap W_i\cap W_j$, with the hypotheses of
Lemma~\ref{lem:ratio}. Hence the ratio
$h_{ij}:=z_1^{(j)}/z_1^{(i)}$ extends smoothly and nonvanishingly
across $D_\alpha$, is $J$-holomorphic on $W_i\cap W_j$, and
$e_j=h_{ij}\,e_i$ on $D_\alpha\cap W_i\cap W_j$ (this is
\eqref{eq:cross}). Restricting to $D_\alpha$: by
Lemma~\ref{lem:strata}, $J_\alpha:=J|_{TD_\alpha}$ is an almost
complex structure on the compact connected manifold $D_\alpha$, and
$h_{ij}|_{D_\alpha}$ is $J_\alpha$-holomorphic: indeed
$\dd(h_{ij}|_{D_\alpha})\circ J_\alpha=(\dd h_{ij}\circ J)|_{TD_\alpha}
=i\,\dd h_{ij}|_{TD_\alpha}=i\,\dd(h_{ij}|_{D_\alpha})$. By
Lemma~\ref{lem:max}, $h_{ij}|_{D_\alpha}$ is constant on every
connected component of $D_\alpha\cap W_i\cap W_j$.

\emph{Step 3: flatness and vanishing Euler class.}
The complex line bundle $\nu_\alpha^*$ thus admits a trivialising
cover (the $e_i$) whose transition functions
$\{h_{ij}|_{D_\alpha}\}$ are constant on each connected overlap
piece. Choose on each $W_i\cap D_\alpha$ the connection whose
covariant derivative kills $e_i$; since on overlaps the frames
differ by constants, these local connections agree on overlaps and
assemble to a global connection $\nabla$ on $\nu_\alpha^*$ with
$\nabla e_i=0$ on $W_i$; its curvature satisfies
$\Omega(e_i)=0$ locally, hence $\Omega\equiv0$: $\nabla$ is flat.
By Chern--Weil theory (the first Chern form of a flat
connection is identically zero),
\[
c_1(\nu_\alpha^*)\ =\ 0\ \in\ H^2(D_\alpha;\RR).
\]
For a complex line bundle, the Euler class of the oriented real
plane bundle equals $c_1$, and duals reverse the sign; hence
\[
e(\nu_\alpha)\ =\ c_1(\nu_\alpha)\ =\ -c_1(\nu_\alpha^*)\ =\ 0 ,
\qquad
\langle e(\nu_\alpha),[D_\alpha]\rangle=0 .
\]
The self-intersection $D_\alpha\cdot D_\alpha:=\langle
e(\nu_\alpha),[D_\alpha]\rangle$ therefore vanishes.
\end{proof}

\begin{proof}[Proof of Corollary~\ref{cor:cp2}]
$(\CP^2,\omega_{\mathrm{FS}})$ is a closed symplectic manifold and
a projective line $L$ is a smooth symplectic hypersurface; SNC for a
single component is automatic. Its normal bundle is
$\mathcal O(1)$ with $c_1$ evaluating to $1$ on $[L]$; so the
conclusion of Theorem~\ref{thm:obstruction} fails: no
$\omega_{\mathrm{FS}}$-tame SNC-adapted $J$ exists.

For the deformation clause: suppose $\{\omega_t\}_{t\in[0,1]}$ is a
path of cohomologous symplectic forms with $\omega_0=
\omega_{\mathrm{FS}}$, along which $L$ remains a symplectic SNC
divisor. Apply Theorem~\ref{thm:obstruction} to
$(\CP^2,\omega_1,L)$: the hypothesis "$J$ is $\omega_1$-tame and
SNC-adapted" implies $L\cdot L=0$; but $L\cdot L$ is the
topological intersection number of the homology class
$[L]\in H_2(\CP^2;\mathbb Z)$ with itself, equal to $1$ regardless
of the symplectic form used (it only depends on the underlying
oriented smooth pair $(\CP^2,L)$): contradiction. Hence no such $J$
exists after any allowed deformation either. For
$D=L_1\cup L_2$ formed by two distinct lines: they meet transversely
at one point, and the standard local Darboux model for two
transversally intersecting symplectic hypersurfaces supplies the SNC
charts at the intersection point, so $D$ is a symplectic SNC divisor
in the sense of the question; each component satisfies
$L_i\cdot L_i=1$, so the same conclusion holds.
\end{proof}

\section{Discussion of sharpness}\label{sec:remarks}

\noindent(1) \emph{What the obstruction uses.} Only the existence of
local $J$-holomorphic defining functions and the compactness of the
component enter: the ratio lemma makes the conormal frames constant
along the component via the maximum principle. Tameness of $J$,
closedness of $M$, and the SNC condition beyond smoothness of the
component are superfluous for the obstruction --- but $(\CP^2,L)$
satisfies all the hypotheses of the question, so the answer to the
question as stated is no.

\noindent(2) \emph{The deformation does not help because the
obstruction is not symplectic.} Theorem~\ref{thm:obstruction} is a
statement about $J$ and $D$ alone (given adaptedness); its
conclusion $D_\alpha\cdot D_\alpha=0$ is a topological fact of the
pair. No change of $\omega$ within its cohomology class (or any
class) affects it, as long as $D_\alpha$ stays compact.

\noindent(3) \emph{Sharpness of the principle.} The obstruction
vanishes when $D_\alpha\cdot D_\alpha=0$; in the product situation
$D=X\times\{0\}\subset X\times\CC$ (or $X\times\Sigma^g$) with the
product form, where $\nu_D$ is trivial, an adapted tame structure
exists whenever $X$ carries a tame almost complex structure $J_X$:
take $J_X\oplus i$ (already noted after Corollary~\ref{cor:cp2}).
Hence the obstruction used here cannot by itself distinguish such
cases; whether the answer is yes whenever all components have
vanishing normal Euler class is a further question not decided here.

\noindent(4) \emph{More counterexamples.} Any closed symplectic
$(M,\omega)$ containing a symplectic SNC divisor with a compact
connected component $D_\alpha$ of nonzero self-intersection is a
counterexample, stably under the allowed deformation; for instance a
line in $\CP^2$ (self-intersection $+1$) or an exceptional sphere in
a symplectic blow-up (self-intersection $-1$). By contrast, a smooth
fibre of a ruled surface has self-intersection $0$, so the present
obstruction does not apply to it.

\end{document}
